%!TEX root = ../main.tex
\subsection*{Dear VLDB Meta-Reviewer and Reviewers,}

\noindent
We sincerely thank the Meta-Reviewer and Reviewers for their thorough reading and insightful feedback.
In response, we have made substantial revisions to the paper, including the introduction of real-world datasets, rigorous scalability analyses, clearer articulation of novelty, and additional baseline comparisons.
The updated contents in the paper are highlighted in \revise{blue}.
Figures and tables in this response letter are numbered with Roman numerals (\eg Table I and Figure II), while those in the main paper use Arabic numerals (\eg Table 1 and Figure 2).

Below, we first summarize the \textbf{main changes}, and then present \textbf{detailed responses} to each reviewer and the meta-reviewer.

\subsection*{Our Main Changes}

\stitle{M\modificationnum\label{modify:new-exps} [Seven Groups of New Experiments]:}
We have significantly revised the experimental sections in response to reviewer feedback, incorporating seven new groups of experiments.

\begin{itemize}[leftmargin=15pt]
	\item[(1)] We have added real-world evaluation on the Buildings~\cite{sharma2023automatic} dataset, a benchmark comprising 105 ADP tasks derived from data processing pipelines collected from 21 energy companies. We compare \sys with strong baselines and report accuracy, completion rate, and inference cost. The experimental results are reported in Table~\ref{\labelprefix tbl:buildings_exp}.
	\item[(2)] We have evaluated the scalability of \sys in terms of data sizes, table numbers and pipeline lengths. We reported the accuracy, complete rate and resource usage in cases with varying data sizes, table numbers and pipeline lengths. The experimental results are shown in Figure~\ref{\labelprefix fig:performance_compare_by_column_size_acc}, Figure~\ref{\labelprefix fig:performance_compare_by_row_size_acc}, Figure~\ref{\labelprefix fig:performance_compare_on_task_diff_claude}, Figure~\ref{\labelprefix fig:performance_compare_by_table_number_acc} and Table~\ref{\labelprefix tbl:avg_time_memory_by_data_characteristics}
    \item[(3)] We have conducted sensitivity analysis of exploration turns. We report the complete rate and accuracy of \sys when varying maximum exploration turns. The experimental results are reported in Figure~\ref{\labelprefix fig:performance_vary_maxturn_accuracy} and Figure~\ref{\labelprefix fig:performance_vary_maxturn_complete_rate}.
    \item[(4)] We have compared \sys with two agentic coding baselines: Claude Code (CC), and the open-source code agent TRAE. The experimental results are reported in Table~\ref{\labelprefix tbl:agentic_coding_system_exp} and Figure~\ref{fig:combined_pareto_front_compare}.
    \item[(5)] We have added an overall ablation study to isolate how much of the performance gain comes specifically from tree-based reasoning framework and progressive agentic training. The experimental results are reported in Table~\ref{\labelprefix tbl:overall_ablation_study}, with detailed analysis provided in Exp-\ref{exp:ablation_analysis_overall} of Section~\ref{subsec:in_depth_analysis}.
    \item[(6)] We have conducted error analysis of \sys in real-world cases to determine its strength and limitations. The analysis is provided in Exp-\ref{exp:building_exp} of Section~\ref{subsec:in_depth_analysis}.
    \item[(7)] We have added a finer-grained analysis of operator coverage and compositional diversity between Synth-Spider and Parrot, measuring unique bigrams and trigrams, concentration ratios, and cross-category ratios. We also conduct a cross-training experiment comparing models trained on Synth-Spider vs. Parrot, both evaluated on the Parrot test set. We provide detailed analysis in Exp-\ref{exp:data_effectiveness} of Section~\ref{subsec:in_depth_analysis}.
\end{itemize} 


% \stitle{M\modificationnum\label{modify:failure-analysis} [Failure Case Analysis]:}
% We have conducted a dedicated failure case analysis on the Buildings benchmark, identifying two key challenges largely absent in synthesized benchmarks: ambiguous column semantics (50\% of failures) and underspecified data preparation requirements (33\% of failures). Please see the response to \textbf{R1D5}.

% \stitle{M\modificationnum\label{modify:sensitivity} [Sensitivity and Scalability Analyses]:}
% We have conducted deeper sensitivity and scalability analyses along four dimensions: (1) maximum exploration turns (with both Claude-4 and Qwen3-14B), (2) data size, (3) number of source tables, and (4) pipeline length. We compare \sys with ReAct and Code Agent on Claude-4 across all dimensions. We also report runtime and CPU memory usage under varying workload sizes. Please see the responses to \textbf{R1D3}, \textbf{R2W4}/\textbf{R2D4}, and \textbf{R3W3}/\textbf{R3D1}.

\stitle{M\modificationnum\label{modify:technical-details} [Technical Details]:} We have revised Section~\ref{sec:experiment} to provide more comprehensive technical details. In addition, we have included detailed LLM prompts and hyper-parameters of all baselines implemented in this paper in our technical report~\cite{technical_report} due to space limitations. We have also made the following clarifications: (i) We clarified that the environment is deterministic and replaced the probabilistic notation with a deterministic transition function. (ii) We revised Section~\ref{sec:tree_based_agentic_reasoning} to explain how failure branches are preserved with error traces and reused to inform future exploration. (iii) We added detailed discussions of tree-search control, implicit pruning mechanisms, and complexity analysis ($O(T \times L)$) in Section~\ref{subsec:agentic_tree_based_definition}.

% We have provided more technical details in the revised appendix of our technical report~\cite{technical_report}, including prompts used by \sys and other baselines, and the hyperparameters used in our experiments. We have also clarified that the environment is deterministic and replaced the probabilistic notation with a deterministic transition function where appropriate. Please see the responses to \textbf{R1W2}/\textbf{R1D1}, \textbf{R2W3}/\textbf{R2D3}, and \textbf{R3W2}/\textbf{R3D3}.


\stitle{M\modificationnum\label{modify:novelty} [Clarification on Technical Novelty]:} 
We have clarified the technical innovations of \sys beyond existing agentic LLM frameworks from two aspects: (i) The agentic reasoning tree with materialized execution states and failure-aware branching, which differs from ToT that operates over text and MCTS that uses scalar reward signals. (ii) Progressive agentic training with environment token masking and hybrid reward design. And reversible noise injection with verification-in-the-loop for data synthesis.
% Please see the responses to \textbf{R2W1}/\textbf{R2D1} and \textbf{R3W1}.

% \stitle{M\modificationnum\label{modify:ablation} [Overall Ablation Study]:}
% We have added an overall ablation study to isolate how much of the performance gain comes specifically from the tree-based reasoning framework (TAR) and progressive agentic training (PAT). Removing PAT leads to severe degradation (e.g., $-$43.78 accuracy on Synth-Spider), and further ablating TAR causes catastrophic collapse (e.g., $-$61.31 accuracy on Synth-Spider). Please see the response to \textbf{R3D4}.

% \stitle{M\modificationnum\label{modify:failure-branch} [Clarification on Failure Branch Reuse]:}
% We have revised Section~\ref{sec:tree_based_agentic_reasoning} to better explain how failure branches are preserved with their error traces and reused to inform future exploration, rather than being discarded. Please see the response to \textbf{R1D2}.
